\section{Control and Observer Methods}

\begin{frame}{Adaptive Control Method}
    \scriptsize
    考虑二阶通道
    \[
    \ddot q = f(q,\dot q) + b u
    \]
    若 $f$ 可线性参数化为 $f = Y(q,\dot q,\dot q_d,\ddot q_d)\Theta$，定义跟踪误差
    \[
    e=q-q_d,\qquad s=\dot e+\lambda e
    \]
    一种典型自适应控制律为
    \[
    u=\frac{1}{b}\left(-Y\hat\Theta-k_s s+\ddot q_d-\lambda \dot e\right)
    \]
    参数更新律取为
    \[
    \dot{\hat\Theta}=\Gamma Y^T s
    \]
    \begin{itemize}
        \setlength{\itemsep}{0.3em}
        \item $Y$: 回归矩阵，$\Theta$: 未知参数，$\hat\Theta$: 参数估计
        \item $\Gamma \succ 0$: 自适应增益，$k_s>0$: 误差收敛增益
        \item 设计步骤: 建模参数化 $\rightarrow$ 定义滤波误差 $s$ $\rightarrow$ 构造控制律 $\rightarrow$ 用 Lyapunov 设计更新律
        \item 在 RWIP 中，$\Theta$ 可包含惯量、质心偏差、等效重力项等不确定参数
    \end{itemize}
\end{frame}

\begin{frame}{Sliding Mode Control Method}
    \scriptsize
    对于同样的跟踪任务，先定义滑模面
    \[
    s=\dot e+\lambda e
    \]
    将控制律写成“等效项 + 切换项”
    \[
    u=u_{eq}+u_{sw}
    \]
    其中常见设计为
    \[
    u_{eq}=\frac{1}{b}\left(-\hat f+\ddot q_d-\lambda \dot e\right),\qquad
    u_{sw}=-\frac{k}{b}\,\mathrm{sat}\!\left(\frac{s}{\phi}\right)
    \]
    其中
    \[
    \mathrm{sign}(s)=
    \begin{cases}
    1, & s>0\\
    0, & s=0\\
    -1, & s<0
    \end{cases}
    \qquad
    \mathrm{sat}\!\left(\frac{s}{\phi}\right)=
    \begin{cases}
    1, & s>\phi\\
    \dfrac{s}{\phi}, & |s|\le \phi\\
    -1, & s<-\phi
    \end{cases}
    \]
    \begin{itemize}
        \setlength{\itemsep}{0.3em}
        \item $u_{eq}$ 负责名义跟踪，$u_{sw}$ 负责压制有界扰动和建模误差
        \item 若写成 $u_{sw}=-\frac{k}{b}\mathrm{sign}(s)$，鲁棒性更强，但抖振更明显
        \item $\phi>0$ 为边界层厚度，$\mathrm{sat}(s/\phi)$ 在边界层内用线性段替代切换项
        \item 设计步骤: 定义滑模面 $\rightarrow$ 估计名义模型 $\rightarrow$ 选切换增益 $k$ 覆盖扰动上界
        \item 在 RWIP 中，滑模项特别适合处理摩擦、外部冲击和未建模耦合项
    \end{itemize}
\end{frame}

\begin{frame}{Observer Method}
    \scriptsize
    若并非所有状态可直接测量，可先设计观测器，再将估计值送入控制器。\\[0.4em]
    对于线性形式
    \[
    \dot x = A x + B u,\qquad y = C x
    \]
    Luenberger 观测器写为
    \[
    \dot{\hat x}=A\hat x+Bu+L(y-C\hat x)
    \]
    若还希望估计总扰动 $d$，可扩展为 ESO 结构
    \[
    \dot{\hat x}=A\hat x+Bu+E\hat d+L_x(y-C\hat x),\qquad
    \dot{\hat d}=L_d(y-C\hat x)
    \]
    \begin{itemize}
        \setlength{\itemsep}{0.3em}
        \item $L,L_x,L_d$ 通过极点配置或带宽法选择
        \item 设计步骤: 选观测状态 $\rightarrow$ 建立误差动态 $\rightarrow$ 配置观测器极点
        \item 在 RWIP 中可估计角速度、扰动力矩、未测状态，再接入自适应控制或滑模控制
        \item 工程上常见组合: Observer + Adaptive, Observer + SMC, ESO + feedback linearization
    \end{itemize}
\end{frame}
